\documentclass{ximera}

\begin{document}

    \title{Homework 4}
    
    \begin{abstract} Solve the following problems.
    \\
    Due Date: Wednesday, October 28th, 11:59 pm
    \end{abstract}
    \maketitle

\begin{exercise}
[10.1] Determine the eigenvalues, determinant, and singular values
of a Householder reflector. Give both a geometric argument and
an algebraic proof for the eigenvalues.
\end{exercise}

\begin{exercise}
[10.3] Let
\[
Z =
\begin{bmatrix}
1 & 2 & 3 \\
4 & 5 & 6 \\
7 & 8 & 7 \\
4 & 2 & 3 \\
4 & 2 & 2
\end{bmatrix}.
\]
Compute the reduced QR factorization of \(Z\) using modified
Gram--Schmidt, Householder reflections, and MATLAB's or Python's built-in
QR routine. Compare the results.
\end{exercise}

\begin{exercise}
[10.4] Consider the orthogonal matrices
\[
F =
\begin{bmatrix}
-c & s \\
s & c
\end{bmatrix},
\qquad
J =
\begin{bmatrix}
c & s \\
-s & c
\end{bmatrix},
\]
where \(s=\sin\theta\) and \(c=\cos\theta\).  Note that $F$ represents a special case of a Householder reflection in dimension 2, and $J$ is a (Givens) rotation matrix.

\begin{enumerate}
    \item Describe the geometric effects of left multiplication
    by \(F\) and \(J\) on the plane $\mathbb{R}^2$.
    \item \textbf{(GRAD STUDENTS ONLY)} Describe a QR factorization algorithm based on Givens
    rotations.
    \item \textbf{(GRAD STUDENTS ONLY)} Show that the algorithm requires six flops per entry
    operated on, compared with four for Householder
    triangularization.
\end{enumerate}
\end{exercise}

\begin{exercise}
[11.1] Suppose \(A\in\mathbb{C}^{m\times n}\) has the block form
\[
A =
\begin{bmatrix}
A_1 \\
A_2
\end{bmatrix},
\]
where \(A_1\in\mathbb{C}^{n\times n}\) is nonsingular and
\(A_2\in\mathbb{C}^{(m-n)\times n}\) is arbitrary.
Prove that
\[
\|A^+\|_2 \leq \|A_1^{-1}\|_2.
\]
\end{exercise}

\begin{exercise}
[11.3] Let \(m=50\) and \(n=12\). Construct a grid \(t\) of
\(m\) equally spaced points on \([0,1]\), and form the
Vandermonde matrix \(A\) for polynomial fitting of degree
\(n-1\). Let \(b=\cos(4t)\).

Compute the least-squares coefficient vector \(x\) to
sixteen-digit precision using:
\begin{enumerate}
    \item The normal equations using Matlab's \verb|\ | command (or Python's \verb|scipy.linalg.lstsq| command).
    \item Modified Gram--Schmidt QR factorization.
    \item Householder QR factorization.
    \item MATLAB's built-in \texttt{qr} (or Python's equivalent, \verb|scipy.linalg.qr| with option \verb|pivotting=true|).
    \item MATLAB's backslash operator (or Python's equivalent, \verb|scipy.linalg.lstsq|).
    \item The singular value decomposition (SVD).
\end{enumerate}

\textbf{GRAD STUDENTS ONLY} Compare the coefficients and identify digits affected by
rounding error. Comment on the stability of the normal equations.
\end{exercise}




\end{document}